\chapter{Nephrology}

\medterm{Acid-Base Disorder}-- Disturbance of the body's acid-base balance, classified by primary disturbance and compensation. Respiratory acidosis: hypoventilation, high PaCO2, low pH. Respiratory alkalosis: hyperventilation, low PaCO2, high pH. Metabolic acidosis: low HCO3, low pH; with high AG (DKA, lactic acidosis, renal failure, toxins) or normal AG (diarrhea, RTA). Metabolic alkalosis: high HCO3, high pH; from vomiting, diuretics, or hyperaldosteronism. Mixed disorders: two or more primary disturbances simultaneously. Diagnosis: ABG + electrolytes, using Winters formula for appropriate compensation. Management: treat the underlying cause.

\medterm{Acute Interstitial Nephritis}-- Inflammatory infiltrate in the renal interstitium and tubules most commonly caused by drug reactions including NSAIDs, antibiotics particularly penicillins and cephalosporins, proton pump inhibitors, and diuretics. Also caused by infections and autoimmune diseases. Presents with AKI, low-grade fever, rash, and eosinophilia though the classic triad is present in less than one third of cases. Urinalysis shows white blood cell casts, sterile pyuria, and eosinophiluria. Diagnosis requires renal biopsy showing interstitial inflammation with eosinophils. Management includes discontinuation of the offending drug and supportive care; corticosteroids may accelerate recovery in severe cases.

\medterm{Acute Kidney Injury}-- Rapid decline in renal function over hours to days. KDIGO criteria define AKI as creatinine increase by 0.3 mg/dL within 48 hours, 1.5 times baseline within 7 days, or urine output less than 0.5 mL/kg/hr for 6 hours. Classified into stages 1 through 3. Causes include prerenal, intrinsic renal, and postrenal etiologies. Early recognition and treatment of reversible causes is essential to prevent permanent kidney damage.

\medterm{Acute Tubular Necrosis}-- Most common cause of intrinsic AKI resulting from ischemic or nephrotoxic injury to tubular epithelial cells. Ischemic ATN develops from prolonged prerenal states. Nephrotoxic ATN results from direct injury by aminoglycosides, amphotericin B, cisplatin, contrast, and myoglobin. Histology shows tubular cell flattening, brush border loss, and intratubular casts. Oliguric phase may be followed by polyuric phase during tubular regeneration. Recovery takes one to three weeks but may be prolonged in severe cases with comorbidities and advanced age.

\medterm{Alkalosis}-- A metabolic (elevated bicarbonate) or respiratory (low PaCO2) disturbance that raises arterial pH above 7.45. Metabolic alkalosis: vomiting/NG suction (HCl loss), diuretics, mineralocorticoid excess, hypokalemia, and alkali ingestion—compensated by hypoventilation; chloride-responsive vs. chloride-resistant distinguishes causes. Respiratory alkalosis: hyperventilation (anxiety, hypoxia, sepsis, PE, salicylates, pregnancy). Consequences: hypokalemia, hypocalcemia symptoms, arrhythmias, and impaired O2 delivery (Bohr). Diagnosis: ABG and electrolytes. Management: correct the underlying cause, replace chloride/potassium, avoid over-correction.

\medterm{Alport Syndrome}-- Hereditary nephritis caused by mutations in type four collagen genes leading to progressive sensorineural hearing loss, ocular anomalies, and renal failure. X-linked form from COL4A5 mutations is the most common. Autosomal recessive and dominant forms exist. Kidney biopsy shows characteristic thinning and lamellation of the glomerular basement membrane on electron microscopy. Hearing loss is bilateral sensorineural and progressive. Ocular findings include anterior lenticonus and retinal flecks. There is no specific cure; management includes ACE inhibitors or angiotensin receptor blockers to reduce proteinuria and slow progression to end-stage renal disease. Hearing aids for hearing loss and genetic counseling are important.

\medterm{ANCA-Associated Vasculitis}-- Small vessel vasculitis with antineutrophil cytoplasmic antibodies causing pauci-immune crescentic GN. GPA features PR3-ANCA with granulomatous inflammation. MPA features MPO-ANCA without granulomas. EGPA features MPO-ANCA with asthma and eosinophilia. Treatment with corticosteroids and rituximab or cyclophosphamide for induction followed by rituximab maintenance.

\medterm{Anti-GBM Disease}-- Autoimmune disease with antibodies against alpha-3 chain of type IV collagen in glomerular and alveolar basement membranes. Presents with crescentic GN and pulmonary hemorrhage (Goodpasture syndrome). Linear IgG deposits on immunofluorescence. Treatment requires emergent plasmapheresis, cyclophosphamide, and corticosteroids. Dramatically improved outcomes compared to near-uniform fatality before these treatments.

\medterm{Anuria}-- Absence of urine output (<100 mL/day). Classification by cause: prerenal (severe hypovolaemia, shock, renal artery occlusion), renal (acute tubular necrosis, acute cortical necrosis, rapidly progressive glomerulonephritis, bilateral renal artery occlusion), and postrenal (bilateral ureteric obstruction, bladder outlet obstruction, urethral obstruction). Anuria is a medical emergency requiring urgent investigation: renal ultrasound to exclude obstruction, bladder catheterisation, urinalysis, serum creatinine, and renal biopsy if indicated. Management: treat underlying cause (fluid resuscitation for prerenal, relief of obstruction for postrenal, dialysis for severe renal failure). Complications: hyperkalaemia, volume overload, uraemia, and metabolic acidosis.

\medterm{Apparent Mineralocorticoid Excess}-- Rare autosomal recessive disorder caused by deficiency of 11-beta-hydroxysteroid dehydrogenase type 2 normally converting cortisol to cortisone in the collecting duct. Cortisol activates mineralocorticoid receptors causing severe hypertension, hypokalemia, metabolic alkalosis, and low renin with low aldosterone. Presents in childhood with growth retardation. Exogenous liquorice and carbenoxolone inhibit the same enzyme causing an acquired form. Treatment is spironolactone and dexamethasone to suppress ACTH and reduce cortisol production.

\medterm{ARBs}-- Block angiotensin II AT1 receptors providing renoprotection similar to ACE inhibitors. Alternative for ACE inhibitor cough. Do not combine ACE inhibitor and ARB due to increased adverse effects without benefit. Similar renoprotective effects to ACE inhibitors.

\medterm{Arteriovenous Fistula}-- Surgical anastomosis between artery and vein for hemodialysis. Preferred access due to lower infection and thrombosis rates. Radiocephalic at wrist or brachiocephalic at elbow. Takes 6 to 8 weeks to mature. Failure to mature occurs in 30\%.

\medterm{Arteriovenous Graft}-- Synthetic conduit connecting artery to vein for dialysis. Used when veins inadequate for fistula. Thrombosis and infection more common than AVF. Can be used earlier than fistula. Declotting with thrombectomy or angioplasty.

\medterm{Atheroembolic Renal Disease}-- Cholesterol crystal embolization to kidneys after vascular procedures. AKI with livedo reticularis and peripheral eosinophilia. Renal biopsy shows cholesterol clefts. Poor prognosis with progressive renal failure.

\medterm{Autosomal Dominant PKD}-- Most common hereditary kidney disease affecting 1 in 400 to 1000. PKD1 chromosome 16 more severe with earlier ESRD than PKD2 chromosome 4. Progressive cyst growth causes hypertension, pain, hematuria, infections, and eventually ESRD. Screen family members with ultrasound. Tolvaptan may slow cyst growth. Transplant may require native nephrectomy if very large.

\medterm{Autosomal Recessive PKD}-- Rare childhood disorder from PKHD1 mutation. Bilateral enlarged kidneys with fusiform dilation of collecting ducts. Associated with congenital hepatic fibrosis. Severe forms cause pulmonary hypoplasia in utero. Variable severity from perinatal lethal to presentation in childhood or adolescence. Management is supportive with treatment of complications.

\medterm{Bartter Syndrome}-- Inherited autosomal recessive disorder of the thick ascending limb of the loop of Henle causing salt wasting, hypokalemic metabolic alkalosis, and hypercalciuria. Results from mutations in genes encoding the NKCC2 transporter or its regulatory subunits. Presents in childhood with failure to thrive, polyuria, and polydipsia. Distinguished from Gitelman syndrome by normal or elevated urinary calcium. Complications include nephrocalcinosis and chronic kidney disease. Treatment includes indomethacin to reduce prostaglandin-mediated salt wasting, and potassium and magnesium supplementation as needed.

\medterm{Benign Prostatic Hyperplasia (BPH)}-- Benign prostatic hyperplasia is a noncancerous enlargement of the prostate gland that commonly affects aging men, compressing the urethra and obstructing urine flow. Symptoms include urinary frequency, urgency, hesitancy, weak stream, and nocturia. Management ranges from watchful waiting and alpha-blockers or 5-alpha-reductase inhibitors to minimally invasive procedures and surgery.

\medterm{Berger Disease}-- Historical name for IgA nephropathy named after Jean Berger who first characterized mesangial IgA deposition in 1968. Characterized by recurrent gross hematuria triggered by mucosal infections with intervals of microscopic hematuria and proteinuria. Generally benign course in most patients but subset progresses to CKD. Complement activation through alternative and lectin pathways plays important pathogenic role leading to investigation of complement-targeted therapies.

\medterm{BK Virus Nephropathy}-- Polyomavirus infection in transplant recipients. Viral cytopathic changes on biopsy. Reduce immunosuppression as primary treatment. Can lead to graft loss. Monitor viral loads in blood for early detection.

\medterm{Blood Urea Nitrogen}-- BUN: serum urea, the end product of protein catabolism, filtered and partially reabsorbed by the tubules. Used with creatinine to assess renal function and volume status. BUN-to-creatinine ratio: high ratio (>20) suggests prerenal azotemia (dehydration, heart failure, GI bleed—blood is protein load) or catabolic states; low ratio suggests intrinsic renal disease, liver disease, or low protein intake. BUN rises with impaired GFR; interpreted with creatinine because urea depends on protein intake and hydration. Not specific to renal disease alone.

\medterm{Bowman Capsule}-- Double-walled epithelial cup surrounding the glomerular capillaries forming the beginning of the nephron. The visceral layer is composed of podocytes whose foot processes interdigitate to form filtration slits. The parietal layer is composed of simple squamous epithelium. Between the layers is Bowman space which collects the glomerular ultrafiltrate. The urinary pole of Bowman capsule is continuous with the proximal tubule. Injury to the podocyte layer results in proteinuria, while damage to the epithelial cells of the parietal layer can contribute to crescent formation in rapidly progressive glomerulonephritis.

\medterm{Calciphylaxis}-- Calcific uremic arteriolopathy causing painful skin necrosis in CKD. Risk factors include warfarin use, elevated calcium-phosphorus product, and obesity. High mortality rate. Treat with sodium thiosulfate, wound care, and discontinuation of warfarin. Avoid calcium-based binders.

\medterm{Calcitriol and Analogues}-- Active vitamin D forms suppressing PTH. Calcitriol, paricalcitol, doxercalciferol. Monitor calcium and phosphorus to avoid hypercalcemia and hyperphosphatemia. Part of CKD-MBD management.

\medterm{Calcium Stones}-- Most common kidney stone type composed of calcium oxalate or calcium phosphate. Risk factors include hypercalciuria, hyperoxaluria, hypocitraturia, and low urine volume. Treat with thiazides for hypercalciuria, dietary modification, and hydration. Potassium citrate for hypocitraturia. Avoid excessive dietary calcium restriction as it increases oxalate absorption.

\medterm{Central Venous Catheter}-- Temporary or tunneled catheter for dialysis. Internal jugular or femoral placement. Tunneled preferred for long-term with lower infection risk. Last resort access with highest complication rates.

\medterm{Cerebral Salt Wasting}-- Renal sodium wasting causing hyponatremia and volume depletion occurring after intracranial surgery or in setting of intracranial disease. Distinguished from SIADH by the presence of volume depletion with negative sodium balance and compensatory elevation of renin and aldosterone. Results from release of natriuretic peptides in response to intracranial disease. Treatment is saline replacement and fludrocortisone to promote sodium retention. Distinguishing from SIADH is critical because fluid restriction is harmful in cerebral salt wasting and may worsen hypovolaemia.

\medterm{Chronic Allograft Nephropathy}-- Progressive transplant kidney function decline over months to years. Interstitial fibrosis and tubular atrophy on histology. Can be immunologic or non-immunologic. Optimize immunosuppression and manage risk factors including hypertension and recurrent disease.

\medterm{Chronic Kidney Disease}-- Progressive loss of kidney function over months to years, defined by GFR <60 mL/min/1.73 m$^2$ or markers of kidney damage for \(\ge\)3 months. Stages 1-5 based on GFR. Leading causes: diabetes mellitus (diabetic nephropathy) and hypertension. Complications: anemia, mineral bone disorder, metabolic acidosis, hyperkalemia, cardiovascular disease, and uremia. Management: blood pressure control (ACEi/ARB), glycemic control, dietary protein restriction, phosphate binders, vitamin D analogs, erythropoiesis-stimulating agents, and preparation for renal replacement therapy (dialysis or transplant).

\medterm{Chronic Kidney Disease (CKD)}-- Chronic kidney disease is the progressive loss of kidney function over months to years, typically from diabetes, hypertension, or glomerular diseases. It is staged by estimated glomerular filtration rate and albuminuria, with end-stage renal disease requiring dialysis or transplantation. Early detection and management of underlying causes, blood pressure control, and dietary modifications slow progression.

\medterm{CKD}-- Abbreviation for Chronic Kidney Disease: a progressive loss of kidney function over months to years, defined by a glomerular filtration rate (GFR) less than 60 mL/min/1.73m$^2$ or markers of kidney damage (albuminuria, haematuria, structural abnormalities) persisting for at least 3 months. CKD is staged by GFR (G1-G5) and albuminuria category (A1-A3). Major causes include diabetes mellitus and hypertension. Complications: anemia, metabolic acidosis, hyperphosphataemia, hypocalcaemia, secondary hyperparathyroidism, hypertension, and cardiovascular disease. Management: blood pressure control (ACEi/ARB), glycaemic control, dietary protein restriction, phosphate binders, erythropoiesis-stimulating agents, and vitamin D analogues. Progression leads to end-stage renal disease requiring renal replacement therapy (dialysis or transplantation).

\medterm{CKD-Mineral Bone Disorder}-- Systemic disorder of mineral and bone metabolism in CKD. Includes calcium, phosphorus, PTH, and vitamin D abnormalities with vascular calcification and bone disease. Renal osteodystrophy encompasses osteitis fibrosa, osteomalacia, adynamic bone disease, and mixed uremic osteodystrophy. Management includes phosphate binders, active vitamin D analogues, and calcimimetics.

\medterm{Collecting Duct}-- Final segment of the nephron where fine-tuning of urine concentration and composition occurs. Principal cells respond to antidiuretic hormone by inserting aquaporin-2 water channels into the luminal membrane, enabling water reabsorption and urine concentration. Aldosterone acts on principal cells to promote sodium reabsorption and potassium secretion through ENaC channels. Intercalated cells regulate acid-base balance through hydrogen ion secretion and bicarbonate reabsorption. The collecting duct is the site of action of aldosterone, vasopressin, and loop diuretics, and is critical in the regulation of blood pressure, electrolyte balance, and acid-base homeostasis.

\medterm{Contrast-Induced Nephropathy}-- AKI following iodinated contrast. Risk factors include CKD, diabetes, dehydration. Prevent with IV saline hydration. Monitor creatinine 48 to 72 hours post-contrast. Hold metformin in high-risk patients.

\medterm{Creatinine}-- A breakdown product of muscle creatine phosphate produced at a constant rate, filtered by the glomerulus and minimally reabsorbed; serum creatinine is the most used marker of kidney function. Elevated levels indicate reduced GFR (renal impairment), but are affected by muscle mass (low in elderly/malnourished), diet, and drugs (cimetidine, trimethoprim block tubular secretion). Used to estimate GFR (CKD-EPI, Cockcroft-Gault) and stage CKD. Rapid rise indicates acute kidney injury; baseline values aid interpretation. Not sensitive to early GFR loss.

\medterm{Cystine Stones}-- Rare stones from inherited cystinuria with defective cystine transport. Hexagonal crystals on urine microscopy. Difficult to dissolve. Treat with aggressive hydration, urine alkalinization to pH greater than 7.5, and thiol agents including tiopronin and D-penicillamine that form soluble complexes with cystine.

\medterm{Cystinuria}-- Autosomal recessive inherited disorder (incidence ~1 in 7,000 to 1 in 15,000) of proximal renal tubular reabsorption of the dibasic amino acids cystine, ornithine, arginine, and lysine (the "COLA" amino acids), caused by mutations in SLC3A1 (type A, autosomal recessive) or SLC7A9 (type B, sometimes autosomal dominant with incomplete penetrance) encoding the rBAT-b$^{0,+}$ amino acid transporter. Excessive urinary excretion of poorly soluble cystine (about 250 mg/L) exceeds its solubility at urinary pH below 7.5, leading to recurrent cystine kidney stones, often beginning in childhood or adolescence. Clinical: recurrent nephrolithiasis, urinary tract obstruction, hematuria, flank pain, renal colic, complicated by renal insufficiency in up to 70 percent of cases. Diagnosis: hexagonal cystine crystals on urinalysis; urinary cystine excretion $> 250$ mg/g creatinine (normal $< 30$); stone analysis (yellow, hexagonal, sulfurous); genetic testing. Treatment: aggressive hydration (urine output $> 2-3$ L/day, target urine specific gravity $< 1.01$); urine alkalinization to pH 7.0-7.5 with potassium citrate (cystine solubility increases $> 10$-fold between pH 7.0 and 8.0); low sodium and moderate protein (avoid methionine-rich foods); thiol agents (tiopronin first-line, D-penicillamine for refractory) forming soluble mixed disulfides with cystine. Impose surgical intervention only for obstruction (PCNL/URETEROSCOPY).

\medterm{Diabetic Nephropathy}-- Kidney disease caused by diabetes mellitus, the leading cause of ESRD worldwide. Pathophysiology: hyperglycemia-induced glomerular hyperfiltration, basement membrane thickening, mesangial expansion, and glomerulosclerosis. Presents with progressive albuminuria (microalbuminuria $\rightarrow$ macroalbuminuria) and declining GFR. Often accompanied by diabetic retinopathy. Screening: annual urine albumin-to-creatinine ratio and serum creatinine in all diabetic patients. Treatment: strict glycemic control, ACE inhibitors or ARBs (renoprotective independent of blood pressure), blood pressure control, and SGLT2 inhibitors for additional renoprotection.

\medterm{Dialysis Disequilibrium Syndrome}-- Neurological complications during hemodialysis in severely uremic patients. Headache, nausea, confusion, seizures, and coma. Results from rapid solute removal causing cerebral edema. Prevent with gradual dialysis initiation and shorter initial sessions.

\medterm{Distal Tubule}-- Segment of the nephron extending from the thick ascending limb of the loop of Henle to the collecting duct. The early distal tubule reabsorbs sodium and chloride through the NCC cotransporter, which is the target of thiazide diuretics. The late distal tubule and connecting tubule contain principal cells that respond to aldosterone for sodium reabsorption and potassium secretion, and intercalated cells that regulate acid-base balance. The macula densa cells at the junction with the thick ascending limb, where it forms the juxtaglomerular apparatus that regulates renin release, glomerular filtration rate, and renal blood flow through tubuloglomerular feedback.

\medterm{DMSA Scan}-- Nuclear medicine scan evaluating renal cortical function and scarring. Gold standard for detecting pyelonephritis and renal scars in children. Shows photopenic areas corresponding to cortical damage. Important for long-term follow-up of childhood UTI.

\medterm{Electrolyte Imbalance}-- Abnormal serum concentration of electrolytes (sodium, potassium, calcium, magnesium, phosphate, chloride). Hyponatremia (<135 mEq/L): most common electrolyte disorder; causes include SIADH, heart failure, hypovolemia, and polydipsia. Hypernatremia (>145 mEq/L): free water deficit from poor intake, diabetes insipidus, or diuretics. Hypokalemia (<3.5 mEq/L): diuretics, vomiting, diarrhea, hyperaldosteronism. Hyperkalemia (>5.0 mEq/L): renal failure, ACEi, K-sparing diuretics. Hypocalcemia, hypercalcemia, hypomagnesemia, and hypermagnesemia each have specific causes and clinical presentations. Management: correct the underlying cause and replace/remove electrolyte as needed, with careful monitoring to avoid rapid shifts.

\medterm{End-Stage Renal Disease}-- The final stage of chronic kidney disease (CKD stage 5), defined by GFR <15 mL/min/1.73 m$^2$ or the need for renal replacement therapy. Uremia develops with accumulation of toxins causing nausea, anorexia, prpruritus, pericarditis, encephalopathy, and bleeding diathesis. Life-sustaining treatments: hemodialysis, peritoneal dialysis, or kidney transplantation. Management also includes anemia correction, phosphate binding, vitamin D analogs, and blood pressure control.

\medterm{Erythropoiesis-Stimulating Agents}-- Epoetin alfa and darbepoetin alfa stimulate red blood cell production in CKD anemia. Target hemoglobin 10 to 11.5 g/dL. Avoid exceeding 13 g/dL due to increased cardiovascular risk. Ensure adequate iron stores before initiating therapy.

\medterm{ESKD}-- End-stage kidney disease defined as GFR less than 15 or need for dialysis. Medicare eligibility at GFR less than 15 or on dialysis. Palliative care option for those not wanting renal replacement therapy. Comprehensive advance care planning essential.

\medterm{Fanconi Syndrome}-- Generalized proximal tubular dysfunction with impaired reabsorption of glucose, amino acids, phosphate, uric acid, and bicarbonate. Causes glycosuria, aminoaciduria, phosphaturia, and metabolic acidosis. Associated with proximal RTA. Causes include myeloma, heavy metal toxicity, Wilson disease, and genetic disorders. Treatment addresses underlying cause and replaces depleted substances.

\medterm{Fibromuscular Dysplasia}-- Non-atherosclerotic arterial disease affecting renal arteries. Causes renovascular hypertension in young women. String-of-beads on angiography from medial fibroplasia. Treated with angioplasty alone which is highly effective. Renal function normalizes after successful procedure.

\medterm{Fluid Overload}-- Excess total body fluid due to impaired renal sodium/water excretion or excessive intake, leading to edema, weight gain, hypertension, and pulmonary congestion. Causes: heart failure, renal failure (acute/chronic), nephrotic syndrome, cirrhosis, iatrogenic (IV fluids), and hypoalbuminemia. Features: peripheral/pitting edema, pulmonary edema (dyspnea, crackles), JVP elevation, ascites, effusions, hypertension. Diagnosis: clinical, BNP, chest X-ray (cardiomegaly, congestion), and assessment of volume status. Management: sodium/fluid restriction, diuretics (loop), treat the underlying cause; dialysis for refractory renal overload.

\medterm{Focal Segmental Glomerulosclerosis}-- Pattern of glomerular injury with sclerosis in some focal glomeruli and segments. Can be primary idiopathic or secondary to HIV, obesity, heroin use, or adaptive responses to nephron loss. Histology shows segmental sclerosis with hyalinosis. Less responsive to steroids than minimal change with approximately 50\% achieving remission. Secondary FSGS treated by addressing underlying cause. Progressive disease leads to ESRD. Recurrence after transplantation occurs in approximately 30\%.

\medterm{Focal Segmental Glomerulosclerosis (FSGS)}-- The most common cause of nephrotic syndrome in African American adults and Hispanic populations. Associated with HIV infection (HIV-associated nephropathy), heroin use, sickle cell disease, massive obesity, interferon therapy, and renal hyperfiltration (solitary kidney). Pathophysiology: podocyte effacement and segmental sclerosis of glomeruli on light microscopy; immunofluorescence is non-specific (IgM/C3 in sclerotic areas). Electron microscopy shows extensive foot process effacement. Clinical features: heavy proteinuria ($>3.5\text{ g/24h}$), edema, hypoalbuminemia, hypertension, and progressive renal insufficiency. Primary FSGS managed with high-dose oral corticosteroids and immunosuppressants (calcineurin inhibitors). High rate of recurrence after renal transplantation (up to 30-50\%).

\medterm{Frequency}-- Increased daytime micturition, typically more than seven episodes per day, distinguished from polyuria (high urine volume) and nocturia (nighttime voiding). Common causes include urinary tract infection, benign prostatic hyperplasia, overactive bladder, and diabetes mellitus. Evaluation comprises urinalysis, post-void residual, and targeted imaging; management is directed at the underlying cause.

\medterm{Gitelman Syndrome}-- Inherited autosomal recessive disorder of the distal convoluted tubule caused by mutations in the NCC sodium-chloride cotransporter gene. Milder than Bartter syndrome typically presenting in adolescence or adulthood. Characterized by hypokalemic metabolic alkalosis, hypomagnesemia, and hypocalciuria distinguishing it from Bartter syndrome. Patients may have muscle cramps, fatigue, and chondrocalcinosis from hypomagnesemia. Treatment includes magnesium and potassium supplementation and NSAIDs or eplerenone for symptom management.

\medterm{Glomerular Filtration Rate}-- Volume of plasma filtered by the glomeruli per unit time normally ninety to one hundred twenty milliliters per minute per one point seven three square meters body surface area. Estimated by the CKD-EPI equation using serum creatinine, age, sex, and race. Key determinants include renal plasma flow, hydrostatic pressure across the glomerular capillary, and the ultrafiltration coefficient reflecting glomerular permeability and surface area. GFR is the single best indicator of kidney function and decreases progressively with age and nephron loss. The slope of eGFR decline over time is a key prognostic indicator in chronic kidney disease progression.

\medterm{Glomerulonephritis}-- Inflammation of the glomeruli, classified by clinical presentation (nephritic vs. nephrotic syndrome) and histopathology (e.g., IgA nephropathy, post-streptococcal, membranoproliferative, crescentic). Presents with hematuria, proteinuria, hypertension, and reduced GFR. Diagnosis: urinalysis, complement levels, autoantibodies (ANA, ANCA, anti-GBM), and renal biopsy. Treatment depends on etiology: immunosuppression for immune-mediated forms, blood pressure control, and supportive care.

\medterm{Glomerulus}-- The tuft of capillaries in the renal corpuscle (Bowman's capsule) where blood is filtered to form urine. See \hyperlink{Anatomy}{anatomy}.

\medterm{Goodpasture Syndrome}-- Anti-GBM disease affecting both kidneys and lungs. Pulmonary hemorrhage with hemoptysis and crescentic GN. Linear IgG deposits on renal biopsy. Emergent plasmapheresis, cyclophosphamide, and steroids. Early treatment essential to prevent irreversible renal and pulmonary damage.

\medterm{Granular Casts}-- Degenerated cellular debris within tubular lumen appearing as coarse or fine-grained structures. Coarse casts represent more acute process, fine casts indicate chronic degenerative process. Common in acute tubular necrosis reflecting tubular epithelial cell damage. Also seen in chronic kidney disease, nephrotic syndrome, and hypertension. Nonspecific finding of tubular injury.

\medterm{Hemodialysis}-- Extracorporeal blood purification using semipermeable membrane. Three sessions per week for 3 to 4 hours. Requires vascular access with AV fistula preferred. Adequacy assessed by Kt/V greater than 1.4. Complications include hypotension, access infection, and dialyzer reactions.

\medterm{Hyaline Casts}-- Transparent colorless cylindrical structures composed primarily of Tamm-Horsfall protein secreted by the thick ascending limb. Most common type of urine cast. Normal in concentrated urine and after exercise. Increased in dehydration, prerenal states, and volume depletion. Generally benign reflecting increased Tamm-Horsfall concentration or mild tubular dysfunction.

\medterm{Hydronephrosis}-- Dilation of renal collecting system proximal to obstruction. Causes include stones, strictures, malignancy, and BPH. Graded on ultrasound from mild to severe. Prolonged obstruction leads to cortical thinning and renal failure. Bilateral more clinically significant. Urgent decompression required for obstructive uropathy with infection.

\medterm{Hyperkalemia}-- See emergency\_medicine.

\medterm{Hyperphosphatemia}-- Serum phosphate above normal range (>4.5 mg/dL), usually due to impaired renal excretion (CKD—most common), hypoparathyroidism, tumor lysis syndrome, rhabdomyolysis, or excessive intake/absorbable laxatives. Consequences: secondary hyperparathyroidism, vascular/metastatic calcification, and when combined with hypocalcemia, tetany. Features are mostly from the associated hypocalcemia. Management: dietary phosphate restriction, phosphate binders (calcium, sevelamer, lanthanum), treat the cause; in severe/tumor lysis, consider dialysis. Chronic control is central in CKD-MBD.

\medterm{Hypertensive Nephropathy}-- Kidney damage caused by chronic hypertension, leading to nephrosclerosis (arteriolar hyalinosis and glomerulosclerosis). Presents with slowly progressive proteinuria (usually subnephrotic), hypertension, and gradual decline in kidney function. More common in African American patients. Diagnosis of exclusion. Prevention and treatment: strict blood pressure control, ACE inhibitors/ARBs are renoprotective.

\medterm{Hypertensive Nephrosclerosis}-- Kidney damage from chronic hypertension. Benign form with hyaline arteriolosclerosis. Malignant form with hyperplastic arteriolosclerosis and onion-skinning. Second leading cause of ESRD after diabetes. Tight BP control with ACE inhibitors or ARBs essential to slow progression.

\medterm{Hypokalemia}-- Serum potassium less than 3.5 mEq/L. Causes include gastrointestinal losses from vomiting and diarrhea, renal losses from diuretics, RTA, and hyperaldosteronism, intracellular shift from alkalosis and insulin, and inadequate intake. Clinical manifestations include muscle weakness, cramps, cardiac arrhythmias with characteristic ECG changes including U waves, flattened T waves, and ST depression, and rhabdomyolysis in severe cases. Treatment depends on severity and symptoms ranging from oral potassium supplementation to IV potassium replacement at rates not exceeding ten milliequivalents per hour through a peripheral line or twenty milliequivalents per hour through central access, with continuous ECG monitoring during IV replacement.

\medterm{IgA Nephropathy}-- The most common primary glomerulonephritis worldwide, characterized by mesangial IgA deposition. Presents with hematuria (macroscopic with mucosal infections, especially upper respiratory), proteinuria, and hypertension. Most common in young adults, with male predominance. Diagnosis: renal biopsy showing mesangial IgA deposits with mesangial proliferation on light microscopy and electron-dense deposits. Prognosis: 30-40\% progress to ESRD over 20 years. Risk factors: hypertension, proteinuria >1g/day, impaired GFR, and biopsy findings of fibrosis. Treatment: BP control with ACEi/ARB; fish oil; corticosteroids for progressive disease; immunosuppression reserved for crescentic IgA.

\medterm{Immunosuppression in Transplant}-- Medications to prevent rejection. Induction with thymoglobulin or basiliximab. Maintenance with tacrolimus, mycophenolate, and prednisone. Monitor drug levels and renal function. Calcineurin inhibitor nephrotoxicity is major concern requiring dose optimization.

\medterm{Intravenous pyelogram}-- Also known as IVP or excretory urography, a radiographic imaging study of the urinary tract following intravenous administration of iodinated contrast material. The contrast is filtered by the glomeruli and excreted through the collecting system, allowing sequential imaging of the renal parenchyma, calyces, renal pelvis, ureters, and bladder. Indications include evaluation of hematuria, suspected ureteral obstruction (stones, strictures, tumors), congenital anomalies of the urinary tract, and pre-operative assessment of ureteral anatomy. CT urography has largely replaced IVP in modern practice due to superior sensitivity for small masses and stones without the need for compression.

\medterm{Intrinsic Renal AKI}-- AKI caused by damage to renal parenchyma. Major categories include acute tubular necrosis from ischemic or nephrotoxic injury, acute interstitial nephritis from drug reactions, glomerulonephritis, and vascular causes including thrombotic microangiopathy. FENa typically greater than 2\% indicating tubular dysfunction. Urine sediment findings help differentiate specific causes: muddy brown granular casts in ATN, WBC casts in AIN, RBC casts in glomerulonephritis.

\medterm{Juxtaglomerular Apparatus}-- Specialized structure located at the vascular pole of the glomerulus where the distal tubule comes into contact with the afferent arteriole. Composed of macula densa cells in the tubular wall, juxtaglomerular cells in the afferent arteriolar wall that produce renin, and extraglomerular mesangial cells. Senses sodium and chloride delivery through the macula densa and renal perfusion pressure through the juxtaglomerular cells. Regulates renin secretion initiating the renin-angiotensin-aldosterone system and thereby blood pressure, fluid volume, and electrolyte homeostasis.

\medterm{Kidney and Urinary Health}-- Kidney and urinary health encompasses the function of the kidneys, ureters, bladder, and urethra in filtering waste, regulating fluid/electrolyte balance, and eliminating urine. Chronic kidney disease (CKD)--defined by reduced glomerular filtration rate (GFR <60 mL/min/1.73 m$^2$) or kidney damage for >3 months--affects approximately 15\% of adults and is most commonly caused by diabetes and hypertension. Urinary tract infections (UTIs) are bacterial infections (usually \emph{Escherichia coli}) of the bladder (cystitis) or kidney (pyelonephritis) presenting with dysuria, frequency, urgency, and flank pain. Diagnostic evaluation includes serum creatinine/eGFR, urinalysis, urine culture, renal ultrasound, and urodynamic studies. Management emphasizes blood pressure control, glucose management, dietary sodium/protein restriction, avoidance of nephrotoxins, and pharmacotherapy (ACE inhibitors, SGLT2 inhibitors for CKD; antibiotics for UTIs).

\medterm{Kidney Failure}-- The inability of the kidneys to maintain fluid, electrolyte, and acid-base homeostasis and to excrete metabolic wastes. Acute kidney injury (AKI): sudden decline in GFR over hours-days—prerenal (hypoperfusion), intrinsic (acute tubular necrosis, glomerulonephritis, interstitial nephritis), postrenal (obstruction). Chronic kidney disease (CKD): progressive loss >3 months from diabetes, hypertension, glomerulonephritis, polycystic disease; staged by GFR (1-5). Features of uremia: fatigue, prpruritus, edema, nausea, anemia, metabolic acidosis, hyperkalemia. Management: treat cause, control BP/glycemia, RAAS blockade, dietary restriction; dialysis or transplantation in stage 5.

\medterm{Kidney stones}-- Also known as nephrolithiasis or renal calculi, crystalline mineral deposits that form within the renal pelvis or collecting system, affecting approximately 10 percent of the population over a lifetime, with a male-to-female ratio of 2:1 and peak incidence in the 4th-6th decades. Stone types and causes: calcium oxalate stones (most common, 60-70 percent, associated with hypercalciuria, hyperoxaluria, hyperuricosuria, hypocitraturia, and low urine volume), calcium phosphate stones (10-15 percent, associated with renal tubular acidosis, hyperparathyroidism, and urinary tract infections with urease-producing organisms), struvite stones (15-20 percent, associated with urease-producing bacteria such as Proteus mirabilis, forming staghorn calculi), uric acid stones (5-10 percent, radiolucent, forming in acidic urine pH less than 5.5, associated with gout, high purine intake, and chronic diarrhea), and cystine stones (1-2 percent, from cystinuria, hexagonal crystals).

\medterm{Kidney Transplantation}-- Best treatment for ESRD with improved survival over dialysis. Living donor superior to deceased donor. Lifelong immunosuppression required. Rejection types include hyperacute, acute, and chronic. Complications include rejection, infection, vascular complications, and recurrent disease.

\medterm{Kimmelstiel-Wilson Nodules}-- Characteristic round eosinophilic nodular mesangial expansions in diabetic nephropathy. Represent advanced mesangial matrix accumulation. Found in peripheral mesangium associated with diffuse mesangial expansion, GBM thickening, and arteriolar hyalinosis. Indicate advanced diabetic nephropathy correlating with heavy proteinuria and declining GFR.

\medterm{Kt/V}-- Dialysis adequacy measure representing the fractional clearance of urea. K is dialyzer clearance, t is treatment time, V is volume of distribution of urea. Target Kt/V greater than 1.4 for thrice-weekly hemodialysis. Urea reduction ratio greater than 70\% is alternative adequacy measure.

\medterm{Liddle Syndrome}-- Autosomal dominant condition caused by gain-of-function mutations in the ENaC sodium channel subunits leading to constitutive sodium reabsorption and potassium secretion in the collecting duct. Presents with severe hypertension, hypokalemia, and metabolic alkalosis occurring in childhood or young adulthood. Aldosterone levels are suppressed reflecting volume expansion. Unlike primary hyperaldosteronism, the hypertension is resistant to spironolactone because the ENaC mutation is constitutively active. Diagnosis is by low aldosterone and low renin levels, which distinguishes it from primary hyperaldosteronism. Treatment is with amiloride or triamterene, which directly block the ENaC channel.

\medterm{Loop of Henle}-- U-shaped tubular segment consisting of a thin descending limb, thin ascending limb, and thick ascending limb. The descending limb is permeable to water but not solutes, allowing water reabsorption into the hypertonic medullary interstitium. The thick ascending limb actively transports sodium, potassium, and chloride through the NKCC2 transporter, which is the target of loop diuretics, and is impermeable to water creating the dilute tubular fluid essential for the countercurrent multiplication system that generates the medullary osmotic gradient essential for urine concentration.

\medterm{MAG3 Renal Scan}-- Nuclear medicine scan evaluating renal function and obstruction. Differential function between kidneys. Non-obstructive dilation does not show delayed excretion. Useful for assessing split renal function and planning surgical interventions.

\medterm{Malignant Hypertension}-- Severely elevated BP greater than 180/120 with end-organ damage. Renal findings include hyperplastic arteriolosclerosis with onion-skinning and fibrinoid necrosis. Treat with IV antihypertensives. Poor prognosis if untreated. May cause AKI, encephalopathy, and retinopathy.

\medterm{May-Thurner Syndrome}-- Compression of left common iliac vein by right common iliac artery causing left lower extremity DVT. Stenting is mainstay of treatment.

\medterm{Medullary Sponge Kidney}-- Congenital malformation characterized by cystic dilatation of the collecting ducts in the renal papillae predisposing to nephrolithiasis and nephrocalcinosis. Presents with recurrent kidney stones, hematuria, and urinary tract infections. Diagnosis is by CT urogram showing dilated collecting ducts in the renal papillae with or without nephrocalcinosis. Usually bilateral. Treatment is preventive with increased fluid intake and thiazide diuretics for calcium stone prevention. Most patients have no reduction in life expectancy.

\medterm{Membranoproliferative Glomerulonephritis}-- Pattern of glomerular injury with mesangial proliferation and capillary wall thickening creating tram-track appearance on silver stain. Hepatitis C is common cause of immune complex MPGN. C3 glomerulopathy represents complement-mediated disease. Associated with monoclonal gammopathies. Treatment addresses underlying cause with antiviral therapy for hepatitis C and immunosuppression for immune-mediated forms.

\medterm{Membranous Nephropathy}-- A common cause of nephrotic syndrome in adults, characterized by immune complex deposition in the subepithelial space of the glomerular basement membrane. Primary (idiopathic, 75\%): associated with anti-PLA2R antibodies (diagnostic). Secondary: infections (hepatitis B, syphilis), autoimmune (SLE), drugs (NSAIDs, penicillamine), and malignancy (lung, colon). Presents with nephrotic-range proteinuria, edema, and normal renal function initially. Diagnosis: renal biopsy with granular IgG staining along capillary loops; anti-PLA2R antibody testing. Treatment: aggressive BP control, RAAS blockade; immunosuppression (cyclophosphamide + steroids, rituximab, or calcineurin inhibitors) for high-risk progressive disease.

\medterm{Microalbuminuria}-- Excretion of 30 to 300 mg of albumin daily undetectable by standard dipstick. Earliest marker of diabetic nephropathy and cardiovascular risk predictor. Detected by urine albumin-to-creatinine ratio on spot specimen. Normal less than 30 mg/g creatinine. Signals early glomerular damage in diabetes warranting RAAS blockade initiation. Screen all diabetics beginning 5 years after type 1 or at diagnosis of type 2 diabetes.

\medterm{Mineralocorticoid Receptor Antagonists}-- Spironolactone and eplerenone block aldosterone receptors. Antifibrotic effects in CKD. Finerenone is non-steroidal MRA with cardiovascular and renal benefits in diabetic CKD. Monitor for hyperkalemia especially when combined with RAAS inhibitors.

\medterm{Minimal Change Disease}-- Most common cause of nephrotic syndrome in children accounting for approximately 80\%. Glomeruli appear normal on light microscopy but podocyte foot processes are diffusely effaced on electron microscopy. Circulating factors from dysregulated T cells may increase glomerular permeability. Presents with sudden onset nephrotic syndrome. Responds excellently to corticosteroids with complete remission in over 90\% of children. Relapse is common but long-term prognosis is excellent.

\medterm{Myeloma Cast Nephropathy}-- Light chain cast formation in tubules causing AKI in multiple myeloma. Friable waxy casts with giant cells. Bence Jones proteinuria. Treat myeloma and consider plasmapheresis to reduce light chain burden.

\medterm{Nephritic Syndrome}-- Characterized by hematuria with dysmorphic RBCs and RBC casts, oliguria, hypertension, and mild proteinuria less than 3.5 g/day. Results from glomerular inflammation increasing permeability to cells and reducing GFR. Hypertension from sodium and water retention. Caused by post-streptococcal GN, IgA nephropathy, lupus nephritis, and anti-GBM disease. Management includes blood pressure control, diuretics, and immunosuppression when indicated.

\medterm{Nephron}-- Functional unit of the kidney comprising the glomerulus, proximal tubule, loop of Henle, distal tubule, and collecting duct. Each kidney contains approximately one million nephrons responsible for filtering blood, reabsorbing solutes and water, and excreting waste to form urine. Nephron loss is irreversible and leads to compensatory hypertrophy of remaining nephrons, which may eventually contribute to progressive kidney disease through glomerular hyperfiltration and hypertension. The juxtamedullary nephrons with long loops of Henle are essential for producing concentrated urine.

\medterm{Nephrotic Range Proteinuria}-- Protein excretion greater than 3.5 g/day or urine protein-to-creatinine ratio greater than 3.5 g/g. Reflects severe glomerular damage with significant protein loss. Degree correlates with severity of injury and risk of complications including thromboembolism, infection, and progression to ESRD. Causes include minimal change disease, FSGS, membranous nephropathy, diabetic nephropathy, and amyloidosis.

\medterm{Nephrotic Syndrome}-- Clinical constellation of heavy proteinuria greater than 3.5 g/day, hypoalbuminemia less than 3.0 g/dL, generalized edema, hyperlipidemia, and lipiduria. Results from increased glomerular permeability to albumin. Reduced oncotic pressure causes fluid extravasation. Hepatic lipoprotein synthesis increases causing hyperlipidemia. Loss of antithrombin III increases thromboembolic risk. Treatment addresses underlying cause with diuretics for edema and RAAS blockade to reduce proteinuria.

\medterm{Nutcracker Syndrome}-- Compression of left renal vein between aorta and SMA causing hematuria and flank pain. Managed conservatively. Stenting or surgery for severe cases.

\medterm{Oliguria}-- Reduced urine output, defined as <400 mL/day or <0.5 mL/kg/hour in adults. Oliguria is an important sign of acute kidney injury (AKI) and prerenal azotaemia. Causes: prerenal (hypovolaemia—hemorrhage, dehydration, sepsis, heart failure, cirrhosis, nephrotic syndrome), renal (acute tubular necrosis—ischemic or nephrotoxic, acute interstitial nephritis, glomerulonephritis, hemolytic uraemic syndrome), and postrenal (obstruction—benign prostatic hyperplasia, stones, tumors, retroperitoneal fibrosis). Assessment: urine output monitoring (hourly in ICU), urinalysis, urine sodium and osmolality, fractional excretion of sodium (FENa <1\% suggests prerenal, >2\% suggests ATN), serum creatinine and BUN, renal ultrasound (rule out obstruction). Management: treat underlying cause; fluid challenge for prerenal causes, diuretics once volume replete, renal replacement therapy for refractory oliguric AKI with hyperkalaemia, acidosis, or uraemia. Anuria (<100 mL/day) suggests complete obstruction, bilateral renal cortical necrosis, or rapidly progressive glomerulonephritis.

\medterm{Osmotic Demyelination Syndrome}-- Demyelinating process affecting the pons and extrapontine regions caused by overly rapid correction of chronic hyponatremia. The brain adapts to chronic hyponatremia by extruding osmoles and when sodium is corrected too quickly water shifts out of brain cells causing osmotic injury to oligodendrocytes. Classically presents with quadriparesis, dysarthria, dysphagia, and locked-in syndrome. Risk increases when sodium correction exceeds ten milliequivalents per liter in twenty-four hours. Prevention is the most important strategy, with a target rate of correction not exceeding eight milliequivalents per litre per day in high-risk patients.

\medterm{Overactive Bladder}-- A clinical syndrome of urinary urgency, usually with frequency and nocturia, with or without urge incontinence, due to detrusor overactivity during filling. Causes: idiopathic (most), neurogenic (stroke, spinal injury, MS, Parkinson), bladder outlet obstruction, and aging; worsened by caffeine and diuretics. Diagnosis: history, urinalysis, voiding diary, urodynamics in refractory cases. Management: bladder training, timed voiding, pelvic floor exercises, antimuscarinics (oxybutynin, tolterodine), beta-3 agonist (mirabegron), and interventional (botulinum toxin, neuromodulation) for refractory disease.

\medterm{Parathyroidectomy}-- Surgical removal of parathyroid glands for refractory secondary hyperparathyroidism unresponsive to medical therapy. Indicated when calcium-phosphorus product persistently elevated despite maximal medical management, symptomatic bone disease, or calciphylaxis. Subtotal or total parathyroidectomy with autotransplantation.

\medterm{Percutaneous Nephrolithotomy}-- Minimally invasive procedure for removal of large renal stones through a nephrostomy tract created in the flank. A tract is created from the skin to the collecting system under fluoroscopic guidance and dilated to allow passage of a nephroscope. Stones are fragmented and removed using ultrasonic lithotripsy, pneumatic lithotripsy, or laser lithotripsy. Preferred for stones greater than two centimeters or staghorn calculi. Success rates exceed ninety percent. Complications include bleeding requiring transfusion, pneumothorax, and infection, which are less common with experienced operators.

\medterm{Peritoneal Dialysis}-- Dialysis using peritoneum as semipermeable membrane. CAPD with four manual exchanges daily or APD using cycler overnight. Advantages include home-based therapy and hemodynamic stability. Complications include peritonitis and catheter infection. Tenckhoff catheter is standard access.

\medterm{Peritonitis in PD}-- Most common complication of peritoneal dialysis. Cloudy effluent, abdominal pain, and fever. Diagnosis by effluent cell count with neutrophils greater than 50/mm3. Treat with intraperitoneal antibiotics. Recurrent peritonitis may require catheter removal.

\medterm{Phosphate Binders}-- Calcium acetate, sevelamer, lanthanum. Bind dietary phosphate in GI tract. Target normal serum phosphorus in CKD-MBD. Avoid calcium-based binders in patients with vascular calcification.

\medterm{Polyarteritis Nodosa}-- Systemic necrotizing vasculitis affecting medium-sized arteries without glomerular involvement or ANCA positivity. Causes aneurysms, stenosis, and occlusion leading to organ ischemia. Renal involvement causes renovascular hypertension. Other manifestations include mononeuritis multiplex and mesenteric ischemia. Angiography shows microaneurysms. Treatment with corticosteroids and cyclophosphamide.

\medterm{Polycystic Kidney Disease}-- Inherited disorder with progressive renal cyst formation. ADPKD most common from PKD1 or PKD2 mutations. Bilateral enlarged kidneys with countless cysts. PKD1 more severe with earlier ESRD. Extrarenal manifestations include hepatic cysts and berry aneurysms. Screen family members. ARPKD is rare childhood disorder from PKHD1 mutation with congenital hepatic fibrosis.

\medterm{Polycystic Liver Disease}-- Dominant polycystic liver disease is an autosomal dominant condition characterized by multiple hepatic cysts that may cause massive hepatomegaly, abdominal pain, and compressive symptoms. Usually associated with autosomal dominant polycystic kidney disease. Liver function is preserved even with extensive cyst involvement. Diagnosis is by imaging showing multiple hepatic cysts. Treatment of symptoms includes cyst fenestration or liver transplantation for severe cases. Somatostatin analogues may reduce cyst volume and symptoms in selected patients.

\medterm{Polyuria}-- Excessive urine output, defined as >3 L/day in adults. Causes are divided into water diuresis and solute diuresis. Water diuresis: central diabetes insipidus (ADH deficiency), nephrogenic diabetes insipidus (ADH resistance from lithium, hypercalcemia, hypokalemia, genetic causes), and primary polydipsia (psychogenic or dipsogenic). Solute diuresis: hyperglycemia (diabetes mellitus, the most common cause), excess salt intake (IV saline, tube feeding), urea diuresis (high-protein feeding, post-obstruction), and mannitol. Evaluation: 24-hour urine volume, serum and urine osmolality, and water deprivation test. Treatment: address underlying cause; desmopressin (DDAVP) for central DI, thiazides for nephrogenic DI.

\medterm{Post-Streptococcal Glomerulonephritis}-- Immune-mediated glomerulonephritis following infection with nephritogenic strains of group A streptococcus typically 1 to 3 weeks after pharyngitis or 3 to 6 weeks after skin infection. Subepithelial humps on electron microscopy representing immune complex deposits. Presents with nephritic syndrome including hematuria, edema, hypertension, and oliguria. Low complement levels (C3) that normalize within 6 to 8 weeks. Supportive management with most patients recovering completely.

\medterm{Postrenal AKI}-- AKI caused by obstruction of urinary outflow at any level from renal tubules to urethra. Causes include BPH, nephrolithiasis, malignancy, retroperitoneal fibrosis, and neurogenic bladder. Bilateral ureteral obstruction or unilateral in a solitary kidney is required for significant impairment. Early diagnosis with renal ultrasound showing hydronephrosis is essential. Prompt decompression with stenting or nephrostomy can restore renal function if initiated before permanent damage.

\medterm{Potassium Binders}-- Sodium zirconium cyclosilicate and patiromer. Exchange potassium in GI tract for sodium or calcium. Treat chronic hyperkalemia allowing continued RAAS inhibitor use in CKD.

\medterm{Prerenal Azotemia}-- Most common cause of AKI resulting from decreased effective circulating volume. Causes include dehydration, hemorrhage, heart failure, sepsis, and hepatorenal syndrome. Tubular function is intact with urine sodium less than 20 mEq/L and FENa less than 1\%. Urine osmolality greater than 500 reflects avid water reabsorption. BUN-to-creatinine ratio typically greater than 20:1. Usually rapidly reversible with restoration of adequate renal perfusion.

\medterm{Proteinuria}-- Excretion of more than 150 mg of protein in urine daily. Causes include glomerular disease, tubulointerstitial disease, overflow proteinuria from multiple myeloma, and functional proteinuria from fever or exercise. Quantification by 24-hour urine or spot protein-to-creatinine ratio. Significant proteinuria warrants evaluation to identify underlying cause and assess for nephrotic syndrome.

\medterm{Proximal Tubule}-- Segment of the nephron extending from Bowman capsule to the loop of Henle with a convoluted proximal straight tubule and convoluted proximal tubule. Responsible for reabsorbing approximately sixty-five percent of filtered sodium, water, bicarbonate, glucose, amino acids, and phosphate through an extensive brush border. Also secretes organic anions, drugs, and toxins. The high metabolic activity of proximal tubular cells makes them particularly vulnerable to ischemic and toxic injury, which is why proximal tubular injury predominates in ischemic acute tubular necrosis. The proximal tubule is also the primary site for renal drug elimination and plays a key role in drug-induced nephrotoxicity.

\medterm{Pruritus in CKD}-- Severe itching in dialysis patients from uremic toxin accumulation, elevated calcium-phosphorus product, and secondary hyperparathyroidism. Difficult to treat. Include phosphate binders, dialysis optimization, gabapentin, and UV phototherapy.

\medterm{Rapidly Progressive Glomerulonephritis}-- Rapid renal function loss over days to weeks with crescent formation on biopsy. Type I anti-GBM with linear deposits. Type II immune complex with granular deposits. Type III pauci-immune with ANCA vasculitis. Requires aggressive immunosuppression. Without treatment progresses rapidly to ESRD.

\medterm{Rasburicase}-- Recombinant urate oxidase enzyme that converts uric acid to allantoin which is more soluble and easily excreted. Highly effective for prevention and treatment of hyperuricemia in tumor lysis syndrome. More effective than allopurinol for established hyperuricemia. Contraindicated in G6PD deficiency due to risk of hemolytic anemia from oxidative stress. Reduces need for dialysis in tumor lysis syndrome.

\medterm{Renal Anemia}-- Normocytic normochromic anemia in CKD from decreased erythropoietin production. Hemoglobin less than 10 g/dL usually requires treatment with ESAs and iron supplementation. Target hemoglobin 10 to 11.5 g/dL. HIF-PHI prolyl hydroxylase inhibitors are newer oral agents for CKD anemia.

\medterm{Renal Artery Aneurysm}-- Dilatation of renal artery, most common visceral artery aneurysm. Usually asymptomatic. Risk of rupture if greater than 2 cm or symptomatic. Treat with surgical or endovascular repair. Higher rupture risk in pregnant women.

\medterm{Renal Artery Embolism}-- Acute occlusion of renal artery from cardiac emboli typically from atrial fibrillation or mural thrombus. Presents with acute flank pain and hematuria. CT angiography shows filling defect. Treat with anticoagulation and possibly thrombolysis or embolectomy if caught early.

\medterm{Renal Artery Stenosis}-- Narrowing of renal artery causing renovascular hypertension. Atherosclerotic most common in elderly at ostium. Fibromuscular dysplasia in young women causing mid-distal stenosis with string-of-beads appearance. Flash pulmonary edema and declining function on ACE inhibitors indicate intervention. Treatment with angioplasty and stenting for atherosclerotic disease.

\medterm{Renal Biopsy}-- Percutaneous needle biopsy for histological diagnosis. Indicated for unexplained AKI, nephrotic syndrome, hematuria with renal involvement, and transplant dysfunction. Ultrasound-guided for accuracy. Complications include bleeding, arteriovenous fistula, and rarely loss of kidney.

\medterm{Renal Calculi (Kidney Stones)}-- Also known as nephrolithiasis or renal calculi, crystalline mineral deposits that form within the renal pelvis or collecting system, affecting approximately 10 percent of the population over a lifetime, with a male-to-female ratio of 2:1 and peak incidence in the 4th-6th decades. See \hyperlink{Kidney stones}{Kidney Stones}.

\medterm{Renal Colic}-- Severe, paroxysmal flank pain caused by ureteral obstruction, typically by a kidney stone, with radiation to the groin; often the worst pain experienced. Features: sudden onset, colicky waves, nausea/vomiting, hematuria, urinary urgency/frequency; pain location reflects stone level (upper: flank; mid: lower abdomen; lower: suprapubic/groin). Diagnosis: non-contrast CT (gold standard), ultrasound (hydronephrosis), urinalysis (hematuria). Management: analgesia (NSAIDs first-line), hydration, alpha-blockers (tamsulosin) for distal stones, and urologic intervention (shock wave lithotripsy, ureteroscopy) for large (>10 mm), refractory, or complicated stones; treat infection/obstruction urgently.

\medterm{Renal Failure}-- Decline in kidney function leading to accumulation of waste products (azotemia) and electrolyte abnormalities. Acute kidney injury (AKI): rapid decline over hours to days, potentially reversible. Chronic kidney disease (CKD): progressive decline over months to years. Uremic symptoms: nausea, fatigue, prpruritus, confusion, and pericarditis. Management: treat underlying cause, dialysis for severe acute or chronic renal failure, and kidney transplant evaluation.

\medterm{Renal Leukemia Infiltration}-- Kidney involvement by chronic lymphocytic leukemia or other hematological malignancies causing bilateral renomegaly and renal dysfunction. Leukemic cells infiltrate the interstitium and peritubular capillaries. May cause AKI from compression of tubules and vessels or chronic kidney disease. Diagnosis is by renal biopsy showing leukemic infiltrates or non-invasive imaging showing renomegaly. Treatment is directed at the underlying hematological malignancy with chemotherapy.

\medterm{Renal Osteodystrophy}-- Bone disease in CKD including osteitis fibrosa, osteomalacia, adynamic bone disease, and osteoporosis. Results from calcium-phosphate imbalance, secondary hyperparathyroidism, and vitamin D deficiency. Manage with phosphate binders, vitamin D analogues, and parathyroidectomy for refractory hyperparathyroidism.

\medterm{Renal Tubular Acidosis}-- Group of disorders with normal anion gap metabolic acidosis from impaired renal acid excretion or bicarbonate reabsorption. Type 1 distal with impaired H+ secretion and urine pH greater than 5.5. Type 2 proximal with impaired bicarbonate reabsorption. Type 4 hyperkalemic from aldosterone deficiency or resistance. Each type has distinct pathophysiology and treatment approach.

\medterm{Renal Ultrasound}-- Noninvasive imaging evaluating kidney size, structure, and hydronephrosis. First-line imaging for CKD evaluation. Shows cortical thickness, echogenicity, and cysts. Resistive index estimates vascular resistance. Essential for evaluating AKI to exclude obstruction.

\medterm{Rhabdomyolysis and AKI}-- Muscle breakdown releases myoglobin causing tubular obstruction and AKI. Elevated CK and dark urine. Aggressive IV hydration is key prevention and treatment. Alkalinize urine to increase myoglobin solubility.

\medterm{Secondary Hyperparathyroidism}-- Compensatory PTH increase in CKD from decreased phosphate excretion, low calcitriol, and hypocalcemia. Leads to renal osteodystrophy and vascular calcification. Managed with dietary phosphorus restriction, phosphate binders, active vitamin D, and calcimimetics. Refractory cases may require parathyroidectomy.

\medterm{SIADH}-- Syndrome of Inappropriate Antidiuretic Hormone secretion characterized by hyponatremia (dilutional) with inappropriately concentrated urine (>100 mOsm/kg) and elevated urine sodium (>40 mEq/L). Euvolemic hyponatremia. Causes: CNS disorders (meningitis, stroke, tumors), pulmonary disorders (pneumonia, TB, small cell lung cancer), medications (SSRIs, carbamazepine, vincristine), and post-operative state. Diagnosis: euvolemia, hyponatremia, low BUN/uric acid, normal TSH/cortisol. Treatment: fluid restriction (first-line), demeclocycline, vaptans (tolvaptan), and IV hypertonic saline for severe symptomatic hyponatremia with risk of osmotic demyelination.

\medterm{Sodium Bicarbonate Therapy}-- Alkalinizing agent for metabolic acidosis in CKD. Target serum bicarbonate greater than 22 mEq/L. Slows CKD progression and improves muscle wasting. Dose titrated to serum levels.

\medterm{Specific Gravity}-- Ratio of urine density to water density reflecting solute concentration. Normal 1.005 to 1.030. Less precise than osmolality but faster to measure. Elevated in dehydration, glycosuria, proteinuria, and contrast. Decreased in diabetes insipidus and fluid overload. Affected by molecular weight of dissolved substances and influenced by glucose and protein in urine.

\medterm{Staghorn Calculus}-- A large renal stone that occupies the renal pelvis and extends into at least two calyces, creating a cast of the collecting system. Composition: struvite (magnesium ammonium phosphate), associated with urease-producing bacteria (Proteus mirabilis, Klebsiella, Pseudomonas) causing urinary tract infection. More common in women and patients with recurrent UTIs, neurogenic bladder, or urinary diversion. Presents with recurrent UTIs, flank pain, hematuria, and can lead to renal failure if untreated. Diagnosis: CT demonstrates branching stone filling the collecting system. Treatment: PCNL (percutaneous nephrolithotomy) is the gold standard; complete stone clearance essential to prevent recurrence; antibiotics for concurrent infection.

\medterm{Struvite Stones}-- Infection stones composed of magnesium ammonium phosphate. Form in alkaline urine with urease-producing organisms including Proteus and Klebsiella. Staghorn calculi fill renal pelvis almost exclusively from Proteus. Treatment requires complete stone removal with antibiotics and urinary acidification. High recurrence if infection persists.

\medterm{Thrombotic Microangiopathy}-- Endothelial damage with platelet-rich microthrombi causing MAHA with schistocytes, thrombocytopenia, and organ ischemia. TTP from ADAMTS13 deficiency treated with plasmapheresis. HUS from Shiga toxin managed supportively. Atypical HUS from complement dysregulation responds to eculizumab.

\medterm{Tram-Track Appearance}-- Histological finding on silver stain of MPGN representing GBM duplication. Results from mesangial cell interposition into capillary wall with new basement membrane formation. Creates double-contour pattern on Jones methenamine silver stain. Indicates chronic glomerular injury with capillary wall remodeling seen in immune complex-mediated and complement-mediated glomerulonephritis.

\medterm{Uraemia}-- The clinical syndrome resulting from accumulation of urea and other nitrogenous waste products in the blood in advanced chronic kidney disease (CKD stage 4-5) and acute kidney injury. Uraemia is characterized by nausea, vomiting, anorexia, fatigue, prpruritus, insomnia, hiccups, muscle cramps, neuropathy (restless legs, peripheral neuropathy), pericarditis, platelet dysfunction (prolonged bleeding time, easy bruising), metabolic acidosis, hyperkalemia, hyperphosphatemia, and encephalopathy (lethargy, confusion, asterixis, coma). The uraemic syndrome results from accumulation of many organic waste compounds (middle molecules such as beta-2-microglobulin, advanced glycation end-products, p-cresol). Dialysis reverses most symptoms. Uraemic pericarditis is an indication for urgent dialysis. Uraemic fetor is the characteristic ammonia-like odour of the breath.

\medterm{Uremia}-- The clinical syndrome of advanced renal failure resulting from the accumulation of nitrogenous waste products and toxins (urea, creatinine, others), usually with GFR <15-20 mL/min (CKD stage 5). Features: fatigue, anorexia, nausea/vomiting, prpruritus, peripheral neuropathy, restless legs, mental status changes (encephalopathy), pericarditis, bleeding diathesis (platelet dysfunction), and asterixis. Indications for dialysis: uremic symptoms, refractory fluid overload, hyperkalemia, metabolic acidosis, pericarditis, or neuropathy. Managed by renal replacement therapy (hemodialysis, peritoneal dialysis, transplantation).

\medterm{Uremic Encephalopathy}-- Neurological syndrome in severe renal failure. Confusion, asterixis, seizures, coma. Results from uremic toxin accumulation. Diagnosis is clinical in context of advanced renal failure. Emergent dialysis indicated with rapid neurological improvement typically.

\medterm{Uremic Frost}-- Urea crystallization on skin in severe untreated uremia. Rare finding indicating extremely advanced uremia. White crystalline deposits on skin. Indicates immediate need for dialysis. Extremely rare in developed countries due to early detection.

\medterm{Uremic Pericarditis}-- Pericardial inflammation in advanced uremia. Pleuritic chest pain, friction rub, diffuse ST elevation. Indicates need for emergent dialysis. Does not respond to anti-inflammatory therapy. Complications include tamponade. Emergent hemodialysis produces resolution within 2 to 3 weeks.

\medterm{Uremic Platelet Dysfunction}-- Impaired platelet adhesion and aggregation in uremia causing bleeding tendency. Easy bruising, epistaxis, GI bleeding. Improves with dialysis and desmopressin. Cryoprecipitate provides temporary benefit before procedures.

\medterm{Ureterolithiasis}-- The presence of kidney stones (calculi) within the ureter, causing ureteral obstruction and renal colic. Stones form from calcium oxalate/phosphate (most), uric acid, struvite (infection), and cystine. Features: acute severe flank pain radiating to the groin, hematuria, nausea, and vomiting; complications: hydronephrosis, infection (obstructive pyelonephritis—emergency), and renal damage. Diagnosis: non-contrast CT, ultrasound, urinalysis. Management: analgesics (NSAIDs), hydration, alpha-blockers for passage, lithotripsy or ureteroscopy for large/impacted stones; prevention by hydration, diet, and metabolic evaluation for recurrence.

\medterm{Urethritis}-- Inflammation of the urethra, most commonly caused by sexually transmitted infections (Chlamydia trachomatis, Neisseria gonorrhoeae) or non-infectious causes. Presents with dysuria, urethral discharge, and meatal itching. Diagnosis: urine nucleic acid amplification testing (NAAT). Treatment: empiric antibiotics covering gonorrhea and chlamydia; partner treatment essential. Complications: epididymitis, prostatitis, pelvic inflammatory disease.

\medterm{Urgency}-- A sudden, compelling desire to urinate that is difficult to defer—a lower urinary tract symptom. Causes: overactive bladder (detrusor overactivity, often idiopathic), urinary tract infection (urgency + dysuria), bladder outlet obstruction (BPH), neurogenic bladder, interstitial cystitis, and pelvic floor disorders. May be accompanied by frequency, nocturia, and urge incontinence. Evaluation: history, urinalysis, voiding diary, urodynamics if refractory. Management: behavioral (bladder training, pelvic floor therapy), antimuscarinics or beta-3 agonists (mirabegron), treat infection, and assess for outlet obstruction.

\medterm{Uric Acid Nephropathy}-- Acute kidney injury from uric acid crystal precipitation in tubules. Occurs in tumor lysis syndrome and severe hyperuricemia. Prevent with hydration and rasburicase. Alkalinization of urine may be harmful as uric acid precipitates at alkaline pH unlike other stone types.

\medterm{Uric Acid Stones}-- 5 to 10\% of kidney stones. Form in acidic urine with pH less than 5.5 and hyperuricosuria. Risk factors include gout, high purine diet, and chronic diarrhea. Radiolucent but visible on CT. Treat with urine alkalinization with potassium citrate and allopurinol. May dissolve with alkalinization which is unique among stone types.

\medterm{Urine Osmolality}-- Concentration of urine relative to plasma normally 50 to 1200 mOsm/kg. High greater than 500 in prerenal states reflecting intact concentrating ability. Low less than 350 indicates impaired concentrating ability in diabetes insipidus, chronic kidney disease, and excessive fluid intake. More precise than specific gravity. In prerenal azotemia, intact ADH action produces concentrated urine while damaged tubules cannot concentrate effectively.

\medterm{Urine Sediment}-- Microscopic examination of centrifuged urine providing diagnostic information about renal pathology. Normal findings include few RBCs, WBCs, and hyaline casts. Abnormal findings include RBC casts indicating glomerulonephritis, WBC casts indicating pyelonephritis or interstitial nephritis, granular casts in ATN, waxy casts in advanced CKD, and oval fat bodies in nephrotic syndrome. Essential component of AKI and proteinuria evaluation.

\medterm{Urine Sodium}-- Concentration of sodium in urine measured in mEq/L. Low less than 20 in prerenal states reflecting avid tubular reabsorption. High greater than 40 in intrinsic renal disease, diuretic use, and salt-wasting nephropathies. Should be interpreted in context of volume status, medications, and concurrent therapies. Key component in differentiating prerenal from intrinsic causes of AKI.

\medterm{Vitamin D Deficiency}-- See clinical\_nutrition.

\medterm{Water Intoxication}-- A potentially life-threatening condition of hyponatremia (low serum sodium) caused by excessive water intake exceeding the kidneys' maximal free water excretion capacity (~0.7-1 L/hour in healthy adults). Water intoxication dilutes extracellular sodium, causing an osmotic shift of water into cells, particularly dangerous in the brain (cerebral edema). Causes: primary polydipsia (psychogenic polydipsia in psychiatric patients, especially schizophrenia), iatrogenic (excessive hypotonic IV fluids, postoperative, use of hypotonic irrigating solutions during TURP or hysteroscopy $\rightarrow$ TURP syndrome), infant water intoxication (overdiluted formula, forced water ingestion), exercise-associated hyponatremia (marathon runners, endurance athletes who overhydrate with plain water), and MDMA/ecstasy use (increased thirst + SIADH-like effect). Symptoms: mild (nausea, headache, confusion, lethargy), moderate (vomiting, disorientation, muscle cramps), severe (seizures, coma, brainstem herniation, respiratory arrest). Diagnosis: low serum sodium (<135 mEq/L), low serum osmolality (<275 mOsm/kg), and high urine osmolality. Treatment: mild-moderate — water restriction, observation. Severe: hypertonic saline (3\% NaCl, 1-2 mL/kg IV bolus, repeat q10min until symptoms resolve or sodium reaches 120-125 mEq/L). Goal is to raise sodium by 4-6 mEq/L to reverse cerebral edema, avoiding overly rapid correction (>8-10 mEq/L in 24 hours) which risks osmotic demyelination syndrome (central pontine myelinolysis).

\medterm{Waxy Casts}-- Broad refractile cylindrical structures with smooth appearance indicating severe chronic kidney disease. Form from degeneration of granular casts in dilated tubules. Broad diameter reflects formation in dilated collecting ducts. Relatively specific for advanced CKD with significant tubular atrophy. Along with broad casts represent end-stage tubular dysfunction.

\medterm{Wilms Tumor}-- Most common renal malignancy in children. Large abdominal mass. Associated with WAGR, Beckwith-Wiedemann, and hemihypertrophy. Treat with surgery and chemotherapy. Excellent prognosis exceeding 90\%.

\medterm{Estimated Glomerular Filtration Rate (eGFR)}-- An estimate of glomerular filtration derived from serum creatinine and, depending on the validated equation, age and sex. It is used to stage chronic kidney disease and guide medication dosing, but is less reliable when creatinine is changing rapidly or in unusual body composition.

\medterm{Albumin-to-Creatinine Ratio (ACR)}-- The urinary albumin concentration standardized to urinary creatinine concentration, usually measured in a spot urine sample. Persistent albuminuria is a marker of kidney damage and is categorized as A1 (<30 mg/g), A2 (30--300 mg/g), or A3 (>300 mg/g) in commonly used CKD classification.

\medterm{Chronic Kidney Disease Staging}-- Classification of kidney disease by cause, eGFR category, and albuminuria category. G1 through G5 represent eGFR from normal or high to kidney failure; staging is applied when abnormalities persist for at least 3 months and helps estimate progression and complications.

\medterm{Acute Kidney Injury Staging}-- A standardized classification based on changes in serum creatinine and urine output. KDIGO stages 1 through 3 quantify the severity of abrupt kidney-function decline and support prognosis, monitoring, and treatment decisions.

\medterm{Fractional Excretion of Sodium (FENa)}-- The percentage of filtered sodium excreted in urine, calculated from serum and urine sodium and creatinine. It may help distinguish prerenal physiology from tubular injury, although diuretics, CKD, sepsis, and other conditions limit interpretation.

\medterm{Fractional Excretion of Urea (FEUrea)}-- The percentage of filtered urea excreted in urine. It may be used as an adjunct when diuretics make FENa less interpretable, but it is also affected by catabolism, sepsis, gastrointestinal bleeding, and chronic kidney disease.

\medterm{Anion Gap}-- A calculated estimate of unmeasured serum anions, commonly sodium minus chloride and bicarbonate. An increased anion gap supports the presence of unmeasured acids in metabolic acidosis, but interpretation must account for albumin concentration and laboratory method.

\medterm{Hypertonic Saline}-- Intravenous saline with a sodium concentration greater than plasma, commonly 3\% sodium chloride. It is used selectively for severe symptomatic hyponatremia or other urgent indications, with careful monitoring to avoid overly rapid correction and osmotic demyelination.

\medterm{Renin-Angiotensin-Aldosterone System (RAAS)}-- A hormonal system in which renin initiates formation of angiotensin II, promoting vasoconstriction and aldosterone secretion. Aldosterone increases distal sodium reabsorption and potassium and hydrogen-ion secretion; RAAS dysregulation contributes to hypertension, heart failure, and CKD progression.

\medterm{Renal Replacement Therapy (RRT)}-- Treatment that substitutes for inadequate kidney excretory and regulatory function. Modalities include intermittent hemodialysis, continuous kidney replacement therapy, peritoneal dialysis, and kidney transplantation.

\medterm{Continuous Kidney Replacement Therapy (CKRT)}-- A prolonged extracorporeal blood-purification therapy, usually delivered continuously in an intensive-care setting. It provides gradual fluid and solute removal and is useful for hemodynamically unstable patients with severe AKI.

\medterm{Hemofiltration}-- An extracorporeal process that removes plasma water and dissolved solutes by convection across a semipermeable membrane, with replacement fluid used to maintain volume. It is used in continuous and intermittent kidney-replacement modalities.

\medterm{Dialysis Adequacy}-- The degree to which a dialysis prescription removes uremic solutes and controls fluid, electrolyte, and acid-base balance while meeting patient-centered treatment goals. Measures such as Kt/V and urea reduction ratio are used for hemodialysis, but clinical status remains essential.

\medterm{Vascular Access for Hemodialysis}-- A surgically created or catheter-based route that permits high-flow blood removal and return during hemodialysis. Arteriovenous fistulae are generally preferred when feasible, followed by grafts; central venous catheters have higher infection and thrombosis risks.

\medterm{Kidney Transplant Rejection}-- Immune-mediated injury to a transplanted kidney caused by recipient antibodies, T cells, or both. Hyperacute, acute, and chronic rejection differ in timing and pathology and require correlation of graft function, imaging, serology, and often biopsy.

\medterm{Renal Replacement Therapy Indications}-- Urgent indications for dialysis include refractory hyperkalemia, severe metabolic acidosis, toxin exposure by a dialyzable agent, refractory fluid overload or pulmonary edema, and symptomatic uremia. The decision incorporates the patient’s trajectory and overall goals of care.

\medterm{Renal Clearance}-- The volume of plasma from which a substance is completely removed by the kidneys per unit time. Clearance depends on filtration, tubular secretion, and tubular reabsorption; creatinine clearance is an imperfect surrogate for GFR because creatinine is also secreted.

\medterm{Urine Albuminuria}-- Abnormal urinary loss of albumin resulting from increased glomerular permeability or impaired tubular reclamation. Persistent albuminuria is independently associated with CKD progression and cardiovascular risk, even when eGFR is preserved.
\medterm{Autosomal dominant polycystic kidney disease}-- A recognized medical diagnosis or clinical finding.
\medterm{Nephrogenic fibrosing dermopathy}-- A recognized medical diagnosis or clinical finding.
\medterm{Oxalate}-- A recognized medical diagnosis or clinical finding.
\medterm{Page kidney}-- A recognized medical diagnosis or clinical finding.
\medterm{Perirenal extramedullary hematopoiesis}-- A recognized medical diagnosis or clinical finding.
\medterm{Primary adrenal insufficiency (Addison’s disease)}-- A recognized medical diagnosis or clinical finding.
\medterm{Right ureterovesical calculus}-- A recognized medical diagnosis or clinical finding.
\medterm{Uremic stomatitis}-- A recognized medical diagnosis or clinical finding.
\medterm{Angiomyolipoma}-- A benign renal tumor composed of blood vessels, smooth muscle, and fat; it is associated with tuberous sclerosis and may cause retroperitoneal hemorrhage.
\medterm{Acute glomerulonephritis}-- Acute inflammation of the renal glomeruli causing hematuria, proteinuria, reduced glomerular filtration, edema, and often hypertension.
\medterm{Acute obstructive uropathy}-- Sudden impairment of urinary flow caused by obstruction of the urinary tract, potentially producing hydronephrosis, pain, infection, and acute kidney injury.
